\chapter{Procesy Lévy'ego i ich wykładniki charakterystyczne}
\label{ch:levy}
\paragraph{Odsyłacze międzytomowe.} Wnioskowanie o zliczeniach i intensywnościach pojawia się w tomie IV, rozdziale~\ref{IV-ch:glm}; proces Poissona przedstawia tom II, rozdział~\ref{II-ch:poisson-process}.


\section{Dlaczego ruch Browna i proces Poissona są członkami jednej rodziny?}

Ruch Browna ma ciągłe trajektorie, a proces Poissona rośnie
skokowo. Łączą je jednak dwie własności: przyrosty na rozłącznych
przedziałach są niezależne, a rozkład przyrostu zależy tylko od
długości przedziału. Naturalne jest więc pytanie, jak wyglądają
\emph{wszystkie} procesy o tych własnościach.
Odpowiedź prowadzi do procesów Lévy'ego: zawierają one
ruchy Browna z dryfem, procesy złożone Poissona, procesy gamma,
procesy stabilne i liczne modele z nieskończenie wieloma skokami.

\begin{definicja}[Proces Lévy'ego]
\label{def:levy}
Niech $d\ge1$. Proces $X=(X_t)_{t\ge0}$ o wartościach w
$\R^d$, określony na przestrzeni probabilistycznej
$(\Omega,\calF,\Prob)$, jest \emph{procesem Lévy'ego}, jeśli:
\begin{enumerate}
\item $X_0=0$ prawie na pewno;
\item dla $0\le t_0<t_1<\cdots<t_n$ zmienne
$X_{t_1}-X_{t_0},\ldots,X_{t_n}-X_{t_{n-1}}$ są niezależne;
\item dla $s,t\ge0$ zachodzi
$X_{s+t}-X_s\overset{\mathrm d}=X_t$;
\item $X_t\to X_s$ według prawdopodobieństwa, gdy $t\to s$;
\item trajektorie są prawostronnie ciągłe i mają granice lewostronne
(własność \emph{càdlàg}) niemal na pewno.
\end{enumerate}
\end{definicja}

Ostatnie dwa warunki mają odrębne role. Ciągłość stochastyczna
opisuje zachowanie rozkładów przy małych zmianach czasu;
własność càdlàg mówi o wyborze \emph{jednej wspólnej wersji}
trajektorii. Każdy proces z pierwszymi czterema własnościami
ma modyfikację càdlàg; dowód ogólnego twierdzenia o regularizacji
nie jest tutaj potrzebny, gdyż od początku pracujemy z taką wersją.

\begin{lemat}[Brak skoku w ustalonej deterministycznej chwili]
\label{lem:levy-no-fixed-jumps}
Niech $X$ będzie procesem Lévy'ego. Dla każdego ustalonego
$t>0$ zachodzi $\Prob(X_t\ne X_{t-})=0$.
\end{lemat}
\begin{proof}
Dla $n$ tak dużych, że $1/n<t$, różnica
$X_t-X_{t-1/n}$ ma prawo $X_{1/n}$. Z ciągłości
stochastycznej w zerze wynika zbieżność tej różnicy do
zera według prawdopodobieństwa. Jednocześnie z istnienia
granicy lewostronnej wynika zbieżność prawie na pewno do
$X_t-X_{t-}$. Granica według prawdopodobieństwa jest
jednoznaczna, więc $X_t-X_{t-}=0$ prawie na pewno.
\end{proof}

Nie jest to sprzeczne ze skokowością procesu Poissona:
może on mieć skoki w \emph{losowych} chwilach, lecz
prawdopodobieństwo skoku w dowolnej chwili ustalonej z góry
jest równe zero.

\section{Półgrupa splotowa i wykładnik charakterystyczny}

Niech $\mu_t=\Law(X_t)$ i
$\varphi_t(u)=\E e^{i\langle u,X_t\rangle}$,
$u\in\R^d$. Niezależność i stacjonarność przyrostów dają
$\mu_{s+t}=\mu_s*\mu_t$ oraz
$\varphi_{s+t}(u)=\varphi_s(u)\varphi_t(u)$.
Ta ostatnia równość jest odpowiednikiem równania wykładniczego.

\begin{twierdzenie}[Wykładnik charakterystyczny procesu Lévy'ego]
\label{th:levy-exponent}
Dla każdego procesu Lévy'ego $X$ istnieje dokładnie jedna
ciągła funkcja $\Psi:\R^d\to\mathbb C$ z $\Psi(0)=0$
taka, że dla wszystkich $t\ge0$ i $u\in\R^d$ zachodzi
\[
\E e^{i\langle u,X_t\rangle}=e^{t\Psi(u)}.
\tag{10.1}\label{eq:levy-exponent}
\]
Ponadto $\Re\Psi(u)\le0$ i $\Psi(-u)=\overline{\Psi(u)}$.
\end{twierdzenie}
\begin{proof}
Ustalmy $u$. Z ciągłości stochastycznej wynika
$\varphi_t(u)\to1$ przy $t\downarrow0$.
Funkcja $\varphi_t(u)$ nie może być zerem dla żadnego $t>0$:
gdyby $\varphi_t(u)=0$, to z
$\varphi_t(u)=[\varphi_{t/n}(u)]^n$ wynikałoby
$\varphi_{t/n}(u)=0$ dla wszystkich $n$, sprzecznie z
granicą w zerze. Ciągła niezerowa funkcja
$t\mapsto\varphi_t(u)$ na przedziale $[0,\infty)$
ma ciągły logarytm $g_u(t)$ z $g_u(0)=0$.
Ponieważ $e^{g_u(s+t)-g_u(s)-g_u(t)}=1$,
wyrażenie w wykładniku przyjmuje wartości w $2\pi i\Z$;
przez ciągłość i wartość zero dla $s=t=0$ jest wszędzie zerem.
Funkcja $g_u$ jest więc ciągła i addytywna, stąd
$g_u(t)=t g_u(1)$. Kładziemy $\Psi(u)=g_u(1)$.

Dla różnych $u$ wybór logarytmu można uzgodnić ciągle
w otoczeniu zera w czasie: jeżeli $u$ pozostaje w ustalonym
małym otoczeniu $u_0$, to $\varphi_t(u)$ jest bliska
$\varphi_t(u_0)$ jednostajnie dla $t$ ze zwartego przedziału,
dzięki wspólnej ciągłości $(t,u)\mapsto\varphi_t(u)$.
Ciągły dobór logarytmu na takim odcinku daje ciągłość
$u\mapsto g_u(1)=\Psi(u)$; możliwa zmiana gałęzi jest
wielokrotnością $2\pi i$, ale zostaje ustalona przez
warunek $g_u(0)=0$. Równość $\varphi_t(0)=1$ daje
$\Psi(0)=0$. Wreszcie $|\varphi_t(u)|\le1$ i
$\varphi_t(-u)=\overline{\varphi_t(u)}$ implikują
wymienione własności $\Psi$.
Jednoznaczność wynika z ciągłości w $t$ i równości
$e^{t\Psi(u)}=e^{t\widetilde\Psi(u)}$ dla każdego $t$.
\end{proof}

\begin{uwaga}[Nieskończona podzielność]
Dla każdego $n\ge1$ zmienna $X_t$ ma rozkład sumy $n$
niezależnych zmiennych o prawie $\mu_{t/n}$.
Zatem $\mu_t$ jest \emph{nieskończenie podzielna}.
To właśnie klasyfikacja nieskończenie podzielnych praw
stoi za ogólnym wzorem Lévy'ego--Chinczyna.
\end{uwaga}

\begin{przyklad}[Trzy wykładniki]
Dla procesu $X_t=bt+\sigma W_t$ w jednym wymiarze
$\Psi(u)=ibu-\sigma^2u^2/2$.
Dla procesu Poissona o intensywności $\lambda$
$\Psi(u)=\lambda(e^{iu}-1)$.
Dla różnicy niezależnych procesów Poissona
$N_t^+-N_t^-$ o intensywnościach $\lambda_+$ i $\lambda_-$
$\Psi(u)=\lambda_+(e^{iu}-1)+\lambda_-(e^{-iu}-1)$.
Mimo zupełnie innej geometrii trajektorii wszystkie trzy
przykłady mają tę samą wykładniczą zależność od czasu.
\end{przyklad}

\section{Złożony proces Poissona: pełny rachunek}

\begin{definicja}[Złożony proces Poissona]
Niech $N$ będzie procesem Poissona o intensywności
$\lambda>0$, a $(Y_j)_{j\ge1}$ niezależnymi zmiennymi
losowymi o wspólnym prawie $\rho$ na $\R^d\setminus\{0\}$,
niezależnymi od $N$. Proces
$J_t=\sum_{j=1}^{N_t}Y_j$ nazywamy \emph{złożonym
procesem Poissona}. Przyjmujemy, że suma pusta jest zerem.
\end{definicja}

\begin{twierdzenie}[Charakterystyka złożonego procesu Poissona]
\label{th:cpp}
Proces $J$ jest procesem Lévy'ego oraz
\[
\E e^{i\langle u,J_t\rangle}
=\exp\left[t\lambda\left(\int e^{i\langle u,y\rangle}
\,\dd\rho(y)-1\right)\right].
\]
Jeśli $\E\|Y_1\|^2<\infty$, to
$\E J_t=\lambda t\E Y_1$ oraz
$\operatorname{Cov}(J_t)=\lambda t\E(Y_1Y_1^{\mathsf T})$.
\end{twierdzenie}
\begin{proof}
Trajektorie są skokowe i mają na każdym zwartym przedziale
skończenie wiele skoków prawie na pewno.
Niezależność przyrostów wynika ze wspólnej niezależności
liczby i rozmiarów skoków na rozłącznych przedziałach;
stacjonarność z jednorodności zegara Poissona.
Ponieważ $\Prob(N_h>0)=1-e^{-\lambda h}\to0$,
proces jest stochastycznie ciągły.

Warunkując względem $N_t$ i korzystając z niezależności
zmiennych $Y_j$, dla
$\phi(u)=\E e^{i\langle u,Y_1\rangle}$ dostajemy
\[
\E e^{i\langle u,J_t\rangle}
=\sum_{n=0}^\infty e^{-\lambda t}
\frac{(\lambda t)^n}{n!}\phi(u)^n
=e^{\lambda t(\phi(u)-1)}.
\]
Dla momentów najpierw warunkujemy:
$\E(J_t\mid N_t)=N_t\E Y_1$ oraz
$\operatorname{Cov}(J_t\mid N_t)=N_t\operatorname{Cov}(Y_1)$.
Z prawa całkowitej kowariancji wynika
$\lambda t\operatorname{Cov}(Y_1)+
\lambda t(\E Y_1)(\E Y_1)^{\mathsf T}$,
co jest podaną macierzą.
\end{proof}

\section{Miara Lévy'ego i małe skoki}

\begin{definicja}[Miara Lévy'ego]
Miara $\nu$ na $\R^d\setminus\{0\}$ jest
\emph{miarą Lévy'ego}, jeśli
\[
\int_{\R^d\setminus\{0\}}(1\wedge\|y\|^2)\,\dd\nu(y)<\infty.
\tag{10.2}\label{eq:levy-measure}
\]
W szczególności $\nu(\{\|y\|>\varepsilon\})<\infty$
dla każdego $\varepsilon>0$, ale może zachodzić
$\nu(\{0<\|y\|\le1\})=\infty$.
\end{definicja}

Warunek z kwadratem jest subtelny: nie żąda całkowalności
rozmiarów małych skoków. Umożliwia jednak sumowanie
\emph{skompensowanych} małych skoków, ponieważ ich wariancje
są sumowalne. Aby to zobaczyć, skonstruujemy proces z zadanej
miary $\nu$ bez korzystania z nieopracowanej jeszcze teorii
ogólnych losowych miar Poissona.

\begin{lemat}[Suma niezależnych skoków w rozłącznych klasach]
\label{lem:annular}
Niech $\nu$ spełnia \eqref{eq:levy-measure}. Dla dowolnego
zbioru borelowskiego $A$ o $\nu(A)<\infty$, rozłącznego z
pewnym otoczeniem zera, można skonstruować złożony proces
Poissona $J^A$ o wykładniku
$\Psi_A(u)=\int_A(e^{i\langle u,y\rangle}-1)\,\dd\nu(y)$.
Dla rozłącznych zbiorów $A_1,\ldots,A_k$ procesy
$J^{A_1},\ldots,J^{A_k}$ można wybrać niezależnie.
Jeśli $\int_A\|y\|^2\,\dd\nu(y)<\infty$, to
$M_t^A=J_t^A-t\int_Ay\,\dd\nu(y)$ jest
kwadratowo całkowalnym martyngałem i
\[
\E\|M_t^A\|^2=t\int_A\|y\|^2\,\dd\nu(y).
\]
\end{lemat}
\begin{proof}
Przy $\nu(A)>0$ bierzemy intensywność
$\lambda=\nu(A)$ i rozkład skoku
$\rho(B)=\nu(A\cap B)/\nu(A)$, po czym stosujemy
konstrukcję złożonego procesu Poissona. Przy $\nu(A)=0$
wybieramy proces zerowy. Niezależność na skończonej rodzinie
rozłącznych klas uzyskujemy na przestrzeni produktowej.
Własność martyngału dla $M^A$ wynika z niezależności
przyszłych przyrostów i ich średniej zero.
Oszacowanie drugiego momentu otrzymujemy ze wzoru na
kowariancję w twierdzeniu \ref{th:cpp} przez wzięcie śladu.
\end{proof}

\section{Wzór Lévy'ego--Chinczyna i konstrukcja zadanej trójki}

\begin{twierdzenie}[Postać Lévy'ego--Chinczyna]
\label{th:LK}
Proces Lévy'ego $X$ o wartościach w $\R^d$ ma wykładnik
charakterystyczny postaci
\[
\begin{split}
\Psi(u)={}&i\langle b,u\rangle-\frac12u^{\mathsf T}Qu\\
&+\int_{\R^d\setminus\{0\}}
\left(e^{i\langle u,y\rangle}-1
-i\langle u,y\rangle\ind{\{\|y\|\le1\}}\right)\,\dd\nu(y),
\end{split}
\tag{10.3}\label{eq:LK}
\]
gdzie $b\in\R^d$, $Q$ jest symetryczną
nieujemnie określoną macierzą, a $\nu$ jest miarą
Lévy'ego. Przy ustalonej funkcji obcięcia
$y\mapsto y\ind{\{\|y\|\le1\}}$ trójka $(b,Q,\nu)$
jest jednoznaczna. Odwrotnie każda taka trójka wyznacza
prawo procesu Lévy'ego.
\end{twierdzenie}

\begin{uwaga}[Status dowodu klasyfikacyjnego]
Dwa aspekty twierdzenia trzeba rozdzielić. \emph{Kierunek
konstrukcyjny} --- od trójki $(b,Q,\nu)$ do procesu ---
udowodnimy poniżej bez stosowania ogólnej teorii miar losowych.
\emph{Kierunek konieczny} i jednoznaczność trójki
są klasycznym twierdzeniem o reprezentacji funkcji
charakterystycznych praw nieskończenie podzielnych.
Pełny dowód tego wyniku analitycznego wymaga teorii
miar Lévy'ego otrzymywanych z granic małoczasowych
oraz wyodrębnienia części gaussowskiej. Przyjmujemy tutaj
klasyfikację praw nieskończenie podzielnych jako
osobne twierdzenie pomocnicze i nie przedstawiamy poniższej
konstrukcji jako dowodu \emph{obu} kierunków.
Dowód przez miary skoków i dekompozycję Lévy'ego--Itô
będzie naturalnym tematem następnego rozdziału.
\end{uwaga}

\begin{twierdzenie}[Konstrukcja procesu dla dowolnej trójki Lévy'ego]
\label{th:triplet-construction}
Niech $b\in\R^d$, $Q\ge0$ oraz $\nu$ spełnia
\eqref{eq:levy-measure}. Istnieje proces Lévy'ego
$X$, którego wykładnik jest równy prawej stronie
\eqref{eq:LK}. Można go otrzymać jako granicę
jednostajną według prawdopodobieństwa na każdym
przedziale zwartym sum niezależnych procesów skokowych,
ruchu Browna i dryfu.
\end{twierdzenie}
\begin{proof}
\emph{Krok 1: skoki duże.} Niech $A_{\rm big}=\{\|y\|>1\}$.
Miara $\nu(A_{\rm big})$ jest skończona, więc lemat
\ref{lem:annular} daje proces $J^{A_{\rm big}}$. Nie potrzebujemy
żadnego założenia momentowego dla dużych skoków:
w skończonym czasie występuje ich skończenie wiele.

\emph{Krok 2: rozbicie małych skoków.} Dla $k\ge0$ niech
$A_k=\{2^{-(k+1)}<\|y\|\le2^{-k}\}$.
Możemy skonstruować niezależne kompensowane procesy
$M^{A_k}$, ponieważ miary $\nu(A_k)$ są skończone.
Definiujemy sumy $S_t^{(n)}=\sum_{k=0}^n M_t^{A_k}$.
Dla $n>m$, niezależność i średnia zero dają
\[
\E\|S_T^{(n)}-S_T^{(m)}\|^2
=T\sum_{k=m+1}^n\int_{A_k}\|y\|^2\,\dd\nu(y).
\]
Norma procesu $S^{(n)}-S^{(m)}$ jest podmartyngałem;
nierówność maksymalna Dooba w $L^2$ zapewnia
\[
\E\sup_{t\le T}\|S_t^{(n)}-S_t^{(m)}\|^2
\le4T\sum_{k=m+1}^n\int_{A_k}\|y\|^2\,\dd\nu(y).
\tag{10.4}\label{eq:levy-l2-sup}
\]
Z warunku \eqref{eq:levy-measure} prawa strona dąży
do zera przy $m,n\to\infty$. Sumy $S^{(n)}$
są więc Cauchy'ego w $L^2$ normy supremum na
$[0,T]$. Wybierając podciąg z sumowalnymi kwadratami
różnic i stosując Borela--Cantellego otrzymujemy
granicę jednostajną prawie na pewno tego podciągu.
Ponieważ jednostajna granica funkcji càdlàg jest càdlàg,
graniczny proces $S$ ma trajektorie càdlàg.
Wykonując konstrukcję dla $T=1,2,\ldots$ i biorąc podciąg
diagonalny, uzyskujemy jeden proces na $[0,\infty)$.

\emph{Uwaga o rozłączności składników.} Składnik $S$ obejmuje
wszystkie skompensowane klasy $A_k$, także $A_0$.
Nie wolno dodawać procesu $J^{A_0}$ po raz drugi: skutkowałoby
to podwójnym policzeniem skoków klasy $A_0$ oraz błędnym
wykładnikiem charakterystycznym.

\emph{Krok 3: część gaussowska.} Niech $W$ będzie
standardowym $d$-wymiarowym procesem Wienera
niezależnym od wszystkich $J^{A_k}$. Niech $Q^{1/2}$
będzie nieujemnym pierwiastkiem macierzy $Q$ i połóżmy
$X_t=bt+Q^{1/2}W_t+J_t^{A_{\rm big}}+S_t$.
Dla każdej sumy skończonej przyrosty na rozłącznych
przedziałach są niezależne i stacjonarne. Zbieżność w
prawdopodobieństwie ich wektorów zachowuje rozkłady
produktowe i stacjonarność. Dostajemy te same własności
dla $X$. Ciągłość stochastyczna wynika z
$\E\|S_{t+h}-S_t\|^2
=h\int_{\{0<\|y\|\le1\}}\|y\|^2\,\dd\nu(y)$,
ciągłości ruchu Browna oraz
$\Prob(J_{t+h}^{A_{\rm big}}-J_t^{A_{\rm big}}\ne0)\le h\nu(A_{\rm big})$.

\emph{Krok 4: obliczenie wykładnika.} Dla $S^{(n)}$
charakterystyka kompensowanego składnika jest równa
\[
\exp\left[t\sum_{k=0}^n\int_{A_k}
(e^{i\langle u,y\rangle}-1-i\langle u,y\rangle)
\,\dd\nu(y)\right].
\]
Zbieżność całek przy $n\to\infty$ wynika z oszacowania
$|e^{iv}-1-iv|\le v^2/2$ dla rzeczywistego $v$.
Duże skoki mają wykładnik
$t\int_{A_{\rm big}}(e^{i\langle u,y\rangle}-1)\,\dd\nu(y)$,
a część gaussowska $-tu^{\mathsf T}Qu/2$.
Mnożąc funkcje charakterystyczne niezależnych składników
i przechodząc do granicy, otrzymujemy dokładnie
\eqref{eq:LK}. Tym samym konstrukcja jest zakończona.
\end{proof}

\begin{uwaga}[Dlaczego kompensujemy małe, lecz nie duże skoki?]
Dla małych skoków może zachodzić
$\int_{\{\|y\|\le1\}}\|y\|\,\dd\nu(y)=\infty$;
zwykła suma wektorowa skoków nie ma wtedy automatycznego
sensu. Po odjęciu średniego przyrostu każdej skończonej
klasy uzyskujemy rodzinę o sumowalnych \emph{wariancjach}.
Duże skoki występują zaś w skończonej liczbie na każdym
przedziale zwartym i można je sumować bez kompensacji.
\end{uwaga}

\section{Generator procesu Lévy'ego: operator nielokalny}

\begin{twierdzenie}[Generator na funkcjach szybko malejących]
\label{th:levy-generator}
Niech $X$ będzie procesem Lévy'ego o trójce $(b,Q,\nu)$,
a $P_tf(x)=\E f(x+X_t)$. Dla każdej rzeczywistej funkcji
$f$ z przestrzeni Schwartza $\calS(\R^d)$ zachodzi
\[
\lim_{t\downarrow0}\frac{P_tf(x)-f(x)}{t}=Af(x)
\qquad (x\in\R^d),
\]
gdzie
\[
\begin{split}
Af(x)={}&\langle b,\nabla f(x)\rangle
+\frac12\operatorname{tr}(QD^2f(x))\\
&+\int_{\R^d\setminus\{0\}}
\bigl(f(x+y)-f(x)
-\ind{\{\|y\|\le1\}}\langle y,\nabla f(x)\rangle\bigr)
\,\dd\nu(y).
\end{split}
\tag{10.5}\label{eq:levy-generator}
\]
Całka po prawej stronie jest bezwzględnie zbieżna.
\end{twierdzenie}
\begin{proof}
Dla $\|y\|\le1$ rozwinięcie Taylora daje
$|f(x+y)-f(x)-\langle y,\nabla f(x)\rangle|
\le\frac12\|D^2f\|_\infty\|y\|^2$.
Dla $\|y\|>1$ różnica $f(x+y)-f(x)$ jest ograniczona
przez $2\|f\|_\infty$. Zatem całka istnieje.

Przyjmijmy konwencję
$\widehat f(u)=\int_{\R^d}e^{-i\langle u,x\rangle}
 f(x)\,\dd\lambda_d(x)$.
Odwrócenie Fouriera i Fubini dają
\[
P_tf(x)=\frac{1}{(2\pi)^d}\int_{\R^d}
 e^{i\langle u,x\rangle}e^{t\Psi(u)}\widehat f(u)
\,\dd\lambda_d(u).
\]
Z \eqref{eq:LK}, $|e^{iv}-1|\le2$ i
$|e^{iv}-1-iv|\le v^2/2$ wynika
$|\Psi(u)|\le C(1+\|u\|^2)$ dla stałej zależnej
od trójki. Ponieważ $\Re\Psi(u)\le0$,
\[
\left|\frac{e^{t\Psi(u)}-1}{t}\right|
=\left|\Psi(u)\int_0^1e^{st\Psi(u)}\,\dd s\right|
\le|\Psi(u)|.
\]
Funkcja $|\Psi(u)\widehat f(u)|$ jest całkowalna,
więc twierdzenie o zbieżności zdominowanej pozwala
przejść do granicy pod całką Fouriera i otrzymać
$\frac1{(2\pi)^d}\int e^{i\langle u,x\rangle}
\Psi(u)\widehat f(u)\,\dd u$.
Wstawiamy wzór \eqref{eq:LK}. Mnożenie transformaty
przez $iu_j$ odpowiada pochodnej $\partial_j f$, a
przez $-u_i u_j$ pochodnej $\partial_{ij}f$.
Dla składnika skokowego Fubini jest dozwolone dzięki
oszacowaniom z pierwszego akapitu i szybkiemu maleniu
$\widehat f$. Otrzymujemy \eqref{eq:levy-generator}.
\end{proof}

\begin{przyklad}[Dlaczego generator jest nielokalny?]
Dla procesu złożonego Poissona o skokach zawsze równych
$a\in\R$ i intensywności $\lambda$ generator jest równy
$Af(x)=\lambda(f(x+a)-f(x))$.
Wartość $Af(x)$ zależy od $f$ w \emph{innym punkcie} $x+a$;
nie da się jej wyrazić wyłącznie przez skończenie wiele
pochodnych funkcji w punkcie $x$. Właśnie dlatego
równania dla procesów skokowych prowadzą do równań
różniczkowo-całkowych, a nie wyłącznie do PDE.
\end{przyklad}

\section{Mocna własność Markowa procesu Lévy'ego}

\begin{twierdzenie}[Mocna własność Markowa]
\label{th:levy-strong-markov}
Niech $X$ będzie procesem Lévy'ego i niech
$(\calF_t)$ będzie uzupełnioną, prawostronnie ciągłą
filtracją naturalną procesu $X$.
Jeżeli $\tau$ jest czasem zatrzymania ze
$\Prob(\tau<\infty)=1$, to proces
$(X_{\tau+t}-X_\tau)_{t\ge0}$ jest niezależny od
$\calF_\tau$ i ma to samo prawo co $X$.
\end{twierdzenie}
\begin{proof}
Najpierw dla czasu deterministycznego $s$ własność wynika
z niezależności przyrostów i twierdzenia o klasie
monotonicznej: przyszłe przyrosty są niezależne od
$\sigma(X_r:r\le s)$. Aby uwzględnić uzupełnienie i
prawostronną ciągłość filtracji, przybliżamy $s$ z góry
chwilami $s_n\downarrow s$; prawostronna ciągłość
trajektorii i zbieżność zdominowana przenoszą równość
warunkową na $\calF_s=\bigcap_{n}\calF_{s_n}$.

Załóżmy teraz, że $\tau\le K$ i niech
$\tau_n=2^{-n}\lceil2^n\tau\rceil$.
Wtedy $\tau_n$ przyjmuje skończenie wiele wartości,
$\tau_n\downarrow\tau$ oraz $\tau\le\tau_n$.
Dla każdego ograniczonego $\calF_\tau$-mierzalnego $Z$
i każdej ograniczonej ciągłej funkcji
$g:(\R^d)^m\to\R$ niezależność przyrostów w
chwilach deterministycznych daje
\[
\E\Big[Z\,g(X_{\tau_n+t_1}-X_{\tau_n},\ldots,
 X_{\tau_n+t_m}-X_{\tau_n})\Big]
=\E Z\,\E g(X_{t_1},\ldots,X_{t_m})
\]
dla $0\le t_1<\cdots<t_m$, gdyż na
$\{\tau_n=r\}$ zmienna $Z\ind{\{\tau_n=r\}}$
jest $\calF_r$-mierzalna. Przez prawostronną ciągłość
trajektorii $X$ przyrosty po $\tau_n$ zbiegają
prawie na pewno do przyrostów po $\tau$.
Zbieżność zdominowana daje tę samą równość dla $\tau$.
Rozszerzamy ją z ciągłych $g$ na ograniczone borelowskie
funkcje przez klasę monotoniczną. Otrzymujemy niezależność
wszystkich wektorów przyszłych przyrostów od $\calF_\tau$.
Rozkłady skończeniewymiarowe wyznaczają prawo procesu
càdlàg na przestrzeni Skorochoda, więc teza zachodzi
dla całego procesu. Dla nieograniczonego, ale skończonego
prawie na pewno $\tau$ stosujemy argument na zdarzeniach
$\{\tau\le K\}$ i przechodzimy $K\to\infty$.
\end{proof}

\section{Modele: proces gamma i procesy stabilne}

\begin{przyklad}[Proces gamma: nieskończenie wiele skoków, ale skończona wariacja]
Niech $\alpha,\beta>0$ i zdefiniujmy miarę na
$(0,\infty)$ przez
$\dd\nu(x)=\alpha e^{-\beta x}x^{-1}\,\dd x$.
Ponieważ $\int_0^1x\,\dd\nu(x)<\infty$ i
$\nu((1,\infty))<\infty$, można zsumować dodatnie
skoki bez kompensacji i otrzymać niemalejący proces Lévy'ego
$G_t$ z wykładnikiem
$\Psi(u)=\int_0^\infty(e^{iux}-1)\,\dd\nu(x)$.
Różniczkując po $u$ pod całką otrzymujemy
$\Psi'(u)=i\alpha\int_0^\infty e^{-(\beta-iu)x}\,\dd x
=i\alpha/(\beta-iu)$, a $\Psi(0)=0$. Stąd
\[
\Psi(u)=-\alpha\log(1-iu/\beta),\qquad
\E e^{iuG_t}=(1-iu/\beta)^{-\alpha t}.
\]
W konwencji \eqref{eq:LK} jego parametr dryfu wynosi
$b=\int_{(0,1]}x\,\dd\nu(x)$, dzięki czemu
skompensowane małe skoki wraz z dryfem tworzą zwykłą
niemalejącą sumę skoków. Wobec tego $G_t$ ma rozkład gamma z parametrem kształtu
$\alpha t$ i intensywnością $\beta$.
Ponieważ $\nu((0,1))=\infty$, liczba skoków procesu
na dowolnym przedziale $[0,t]$ dla $t>0$ jest prawie
na pewno nieskończona: liczby skoków większych niż
$\varepsilon$ są Poissonowskie o parametrach
$t\nu((\varepsilon,\infty))\to\infty$.
Jednak suma ich rozmiarów jest skończona prawie na pewno,
gdyż oczekiwana suma skoków nie większych niż $1$
ma skończoną wartość $t\int_0^1x\,\dd\nu(x)$.
\end{przyklad}

\begin{przyklad}[Symetryczne procesy stabilne]
Dla $0<\alpha<2$ weźmy w jednym wymiarze
$\dd\nu(y)=c|y|^{-1-\alpha}\,\dd y$, $c>0$,
oraz $b=0$, $Q=0$. Warunek \eqref{eq:levy-measure}
jest spełniony: przy zerze całkujemy $|y|^{1-\alpha}$,
a w nieskończoności $|y|^{-1-\alpha}$.
Symetria znosi urojone części całki w \eqref{eq:LK}, zatem
\[
\Psi(u)=\int_\R(\cos(uy)-1)c|y|^{-1-\alpha}\,\dd y
=-C_\alpha |u|^\alpha,
\quad
C_\alpha=2c\int_0^\infty(1-\cos z)z^{-1-\alpha}\,\dd z>0.
\]
Po podstawieniu $z=|u|y$ otrzymujemy skalowanie
$X_t\overset{\mathrm d}=t^{1/\alpha}X_1$.
Dla $\alpha\ne2$ proces ten nie jest gaussowski, a jego
prawa mają zasadniczo inne zachowanie ogonów niż normalne.
\end{przyklad}

\section{Momenty, wariancje i zadanie z trójką Lévy'ego}

\begin{twierdzenie}[Średnia i kowariancja, jeśli istnieją]
\label{th:levy-moments}
Niech $X$ ma trójkę $(b,Q,\nu)$ w konwencji
\eqref{eq:LK}. Jeśli
$\int_{\{\|y\|>1\}}\|y\|\,\dd\nu(y)<\infty$,
to $\E X_t=t(b+\int_{\{\|y\|>1\}}y\,\dd\nu(y))$.
Jeśli dodatkowo
$\int_{\{\|y\|>1\}}\|y\|^2\,\dd\nu(y)<\infty$,
to
\[
\operatorname{Cov}(X_t)
=t\left(Q+\int_{\R^d\setminus\{0\}}yy^{\mathsf T}\,\dd\nu(y)\right).
\]
\end{twierdzenie}
\begin{proof}
W konstrukcji \ref{th:triplet-construction} małe skoki są
scentrowane. Ich częściowe sumy mają średnią zero i
zbieżność w $L^2$, a więc także w $L^1$.
Wektor dryfu ma średnią $bt$, część gaussowska średnią
zero, natomiast suma dużych skoków ma średnią
$t\int_{\|y\|>1}y\,\dd\nu(y)$ pod pierwszym
założeniem. To daje pierwszy wzór.
Przy drugim założeniu wszystkie trzy losowe składniki
są kwadratowo całkowalne i niezależne. Kowariancja
części gaussowskiej wynosi $tQ$, małych skoków
$t\int_{\|y\|\le1}yy^{\mathsf T}\,\dd\nu(y)$,
a dużych $t\int_{\|y\|>1}yy^{\mathsf T}\,\dd\nu(y)$.
Sumując te macierze otrzymujemy tezę.
Prawo i momenty procesu nie zależą od konkretnej
realizacji użytej w konstrukcji, gdyż charakterystyka
\eqref{eq:LK} wyznacza jednoznacznie jego rozkłady.
\end{proof}

\begin{zadanie}[Skoki dodatnie i ujemne]
Niech $X_t=N_t^+-N_t^-$, gdzie $N^+,N^-$ są niezależnymi
procesami Poissona o intensywnościach $\lambda_+,\lambda_->0$.
Wyznaczyć trójkę Lévy'ego, średnią, wariancję i generator.

\emph{Rozwiązanie.} Miara skoków ma postać
$\nu=\lambda_+\delta_1+\lambda_-\delta_{-1}$.
Przy obcięciu $\ind{\{|y|\le1\}}$ skoki $\pm1$
są traktowane jako małe, zatem
$b=\int_{\{|y|\le1\}}y\,\dd\nu(y)=\lambda_+-\lambda_-$,
$Q=0$. Podstawienie do \eqref{eq:LK} daje
$\Psi(u)=\lambda_+(e^{iu}-1)+\lambda_-(e^{-iu}-1)$.
Z twierdzenia \ref{th:levy-moments}:
$\E X_t=t(\lambda_+-\lambda_-)$ oraz
$\Var(X_t)=t(\lambda_++\lambda_-)$. Wreszcie
$Af(x)=\lambda_+(f(x+1)-f(x))+
\lambda_-(f(x-1)-f(x))$.
To zadanie pokazuje, że parametr $b$ zależy od wybranej
funkcji obcięcia: nie zawsze jest dosłownie „dryfem
trajektorii” w intuicyjnym znaczeniu.
\end{zadanie}

\begin{zadanie}[Jak dobrać kompensator?]
Niech $J_t=\sum_{j=1}^{N_t}Y_j$, gdzie
$N$ jest procesem Poissona o intensywności $\lambda$,
a zmienne $Y_j$ mają średnią $m$ i drugi moment $v_2<\infty$.
Znaleźć martyngał związany z $J$ i jego drugi moment.

\emph{Rozwiązanie.} Przyszłe przyrosty $J$ są
niezależne od przeszłości i mają średnią
$\lambda(t-s)m$, stąd
$M_t=J_t-\lambda tm$ jest martyngałem.
Z twierdzenia \ref{th:cpp} jego wariancja wynosi
$\E M_t^2=\lambda tv_2$.
W szczególności poprawnym kompensatorem jest
$\lambda tm$, a nie $\lambda t\Var(Y_1)$:
drugi moment pojedynczego skoku obejmuje także losowość
jego liczby.
\end{zadanie}

\begin{center}
\begin{tikzpicture}[>=Latex,scale=.93]
\draw[->] (0,0)--(9,0) node[right] {$t$};
\draw[->] (0,-1)--(0,3.1) node[above] {$X_t$};
\draw[blue!65!black,very thick] (0,0)--(1.4,.23)--(1.4,1.35)--
(3.0,1.58)--(3.0,.75)--(4.9,1.05)--(4.9,2.23)--
(6.8,2.52)--(6.8,1.82)--(8.5,2.12);
\draw[orange!75!black,dashed,thick] (0,0)--(8.5,1.34);
\node[anchor=west,scale=.87,blue!65!black] at (5.1,2.86)
{część skokowa z dryfem};
\node[anchor=west,scale=.8,orange!80!black] at (5.8,.42)
{trend średni};
\end{tikzpicture}
\par\smallskip\small Schemat trajektorii procesu złożonego Poissona
z dryfem. Skoki mogą mieć oba znaki; pomarańczowa linia
nie jest realizacją procesu, lecz ilustracją jego trendu.
\end{center}

\section{Granica rozdziału i przygotowanie do całkowania ze skokami}
Konstrukcja z twierdzenia \ref{th:triplet-construction}
pokazuje, jak zbudować proces dla trójki Lévy'ego,
wykorzystując niezależne klasy skoków i kontrolę w $L^2$.
Nie zdefiniowaliśmy jednak jeszcze ogólnych \emph{losowych
miar skoków} ani całek względem ich kompensacji.
Będą one przedmiotem rozdziału 11, gdzie otrzymamy
jedną notację obejmującą zarówno złożone procesy Poissona,
jak i skoki o nieskończonej aktywności.

Równocześnie do całek po trajektoriach wrócimy po omówieniu
pełnej teorii półmartyngałów. Dzięki temu nie będziemy
utożsamiać całki względem miary Wienera z całką Itô
ani traktować oscylacyjnej całki Feynmana jak zwykłej
miary probabilistycznej.

% ===========================================================================
% ROZDZIAL 11 -- Miary skokow, kompensacja, calki Poissona
% ===========================================================================
